\section{Modules}
\begin{frame}{What is a Module?}
    \begin{itemize}
        \item A module is a file containing Python code.
        \item Modules can define functions, classes, and variables.
        \item Modules can also include runnable code.
    \end{itemize}
\end{frame}

\begin{frame}{Namespaces}
    \begin{itemize}
    \item Functions imported as part of a module live in their own namespace. 
    \item A \textbf{namespace} is simply a space within which all names are distinct from each other. 
    \item The same name can be reused in different namespaces but two objects can’t have the same name within a single namespace
    \item One example of a namespace is the set of street names within a single city. Many cities have a street called “Main Street”, but it’s very confusing if two streets in the same city have that name!
    \item Another example is the folder organization of file systems. You can have a file called \texttt{todo} \end{itemize}
\end{frame}


\begin{frame}{Importing Modules in Python}
    \begin{itemize}
        \item If you have a file \texttt{myprog.py} in directory \texttt{~/Desktop/mycode/}, and \texttt{myprog.py} contains \texttt{import morecode}:
        \begin{itemize}
            \item The Python interpreter will look for a file called \texttt{morecode.py}.
            \item It will execute the code in \texttt{morecode.py}.
            \item The object bindings from \texttt{morecode.py} will be available for reference in \texttt{myprog.py}.
        \end{itemize}
        \item Note that the import statement is \texttt{import morecode}, not \texttt{import morecode.py}.
        \item The file must be named \texttt{morecode.py}.
    \end{itemize}
\end{frame}

\begin{frame}{Importing Modules}
    \begin{itemize}
        \item To use a module, you must import it.
        \item You can import an entire module or specific functions from a module.
        \item You can also import a module with an alias.
    \end{itemize}
\end{frame} 


\begin{frame}[fragile]{Importing Modules}
    \begin{itemize}
        \item Import the entire module:
    \end{itemize}

    \begin{lstlisting}[language=Python]
import math
print(math.sqrt(16))
    \end{lstlisting}

    \begin{itemize}
        \item Import specific functions:
    \end{itemize}

    \begin{lstlisting}[language=Python]
from math import sqrt
print(sqrt(16))
    \end{lstlisting}
\end{frame}     


\begin{frame}[fragile]{Importing Modules}
    \begin{itemize}
        \item Import a module with an alias:
    \end{itemize}

    \begin{lstlisting}[language=Python]
import math as m        
print(m.sqrt(16))
    \end{lstlisting}
\end{frame}

\begin{frame}{Standard Library}
    \begin{itemize}
        \item Python has a large standard library.
        \item The standard library is a collection of modules that come with Python.
        \item You can use these modules in your programs.
    \end{itemize}
\end{frame}

\begin{frame}{Standard Library}
    \begin{itemize}
        \item Some common modules in the standard library:
        \begin{itemize}
            \item \texttt{math} - Mathematical functions
            \item \texttt{random} - Random number generators
            \item \texttt{os} - Operating system interfaces
            \item \texttt{sys} - System-specific parameters and functions
            \item \texttt{datetime} - Date and time functions
        \end{itemize}
    \end{itemize}
\end{frame}

\begin{frame}{Third-Party Libraries}
    \begin{itemize}
        \item Python has a large ecosystem of third-party libraries.
        \item Third-party libraries are not part of the standard library.
        \item You can install third-party libraries using \texttt{pip}.
    \end{itemize}
\end{frame} 

\begin{frame}{Third-Party Libraries}
    \begin{itemize}
        \item Some common third-party libraries:
        \begin{itemize}
            \item \texttt{numpy} - Numerical computing
            \item \texttt{pandas} - Data manipulation and analysis
            \item \texttt{matplotlib} - Plotting and visualization
            \item \texttt{requests} - HTTP library
            \item \texttt{flask} - Web development
        \end{itemize}
    \end{itemize}
\end{frame}

\begin{frame}{Summary}
    \begin{itemize}
        \item A module is a file containing Python code.
        \item Modules can define functions, classes, and variables.
        \item Modules can also include runnable code.
        \item Functions imported as part of a module live in their own namespace.
        \item To use a module, you must import it.
        \item You can import an entire module or specific functions from a module.
        \item You can also import a module with an alias.
    \end{itemize}
\end{frame}
